\documentclass[aspectratio=169,professionalfonts]{beamer}
\usetheme{Madrid}
\usecolortheme{default}
\usefonttheme{professionalfonts}
\usepackage[utf8]{inputenc}
\usepackage[T1]{fontenc}
\usepackage{lmodern}
\usepackage{amsmath,amssymb,mathtools,bm}
\usepackage{booktabs}
\usepackage{array}
\usepackage{tabularx}
\usepackage{tikz}
\usetikzlibrary{arrows.meta,positioning,fit,calc}
\usepackage{hyperref}

\title{LDPC Scheduling, Weak Coupling, and State Augmentation}
\subtitle{A step-by-step academic framing}
\author{}
\date{}

\begin{document}

\begin{frame}
  \titlepage
\end{frame}

\begin{frame}{Roadmap}
\begin{enumerate}
  \item Define weakly coupled MDPs, restless bandits, and factored MDPs.
  \item Map clustered LDPC scheduling to these concepts.
  \item Explain why the mapping is not exact.
  \item Introduce the state-augmentation idea and its interpretation.
  \item State the most accurate classification and nearest literature connections.
\end{enumerate}
\end{frame}

\begin{frame}{Step 1: Weakly coupled MDPs (WC-MDP)}
A weakly coupled MDP is built from multiple sub-MDPs whose \emph{local dynamics are independent}, while the system is coupled through shared action/resource constraints.

\vspace{0.4em}
Typical structure:
\[
S = S_1 \times \cdots \times S_K,
\qquad
A(s) = \Bigl\{a=(a_1,\dots,a_K): \sum_{u=1}^K c_u(a_u) \le B\Bigr\},
\]
with transition kernel
\[
P(s'\mid s,a)=\prod_{u=1}^K P_u(s_u'\mid s_u,a_u).
\]
\textbf{Note: }Standard WC-MDPs are coupled by shared resource constraints (e.g. common budget) but transitions are independent. Many theoretical results break for other formulations (e.g. state based transition coupling used here)
\end{frame}

\begin{frame}{Step 2: Restless bandits}
A restless bandit is a special case of a weakly coupled MDP in which each subsystem is an arm, each arm has a small action set (usually active/passive), and only a limited number of arms can be activated at each time.

\vspace{0.4em}
For arm $u$:
\[
a_u \in \{0,1\},
\qquad
\sum_{u=1}^K a_u \le B,
\]
and even passive arms keep evolving:
\[
P_u(s_u'\mid s_u,a_u=0) \neq \delta_{s_u}.
\]
Whittle-index methods rely on this per-arm local Markov structure plus the activation constraint.
\end{frame}

\begin{frame}{Step 3: Factored and dynamically coupled MDPs}
A factored MDP represents the global state by components $s=(s_1,\dots,s_K)$ (here states of clusters) and exploits sparse conditional dependence (the fact that clusters are almost independent).

\vspace{0.4em}
In a \emph{dynamically coupled} factored MDP, the next state of subsystem $u$ may depend on neighboring subsystem states/actions:
\[
P(s_u' \mid s,a) = P_u\bigl(s_u'\mid s_{\mathcal N(u)}, a_{\mathcal N(u)}\bigr),
\]
for a sparse interaction graph $\mathcal N(u)$. This is strictly more general than a classical weakly coupled MDP.
\end{frame}

\begin{frame}{Step 4: Clustered LDPC scheduling model}
In clustered LDPC decoding, the Tanner graph is partitioned into clusters of check nodes. At each decision step, the scheduler picks one cluster to update.

\vspace{0.4em}
Natural abstraction:
\begin{itemize}
  \item subsystem $u$ $\leftrightarrow$ cluster $C_u$,
  \item local state $s_u$ $\leftrightarrow$ cluster summary, e.g. hard decisions / residual-like statistics / local syndrome features,
  \item action $a_u$ $\leftrightarrow$ whether cluster $u$ is scheduled now,
  \item budget constraint $\sum_u a_u \le B$ $\leftrightarrow$ serial or parallel hardware scheduling constraint.
\end{itemize}
This is why the problem initially looks like a weakly coupled MDP or restless bandit.
\end{frame}

\begin{frame}{Step 5: Where the direct WC-MDP mapping fails}
The issue is not the action budget; the issue is the transition law.

\vspace{0.4em}
If cluster $u$ is updated, messages are changed on VNs adjacent to $C_u$. Those VNs may also be adjacent to other clusters $v$, so the next state of cluster $v$ can change because of action $a_u$.

\vspace{0.4em}
Hence one should not expect
\[
P(s'\mid s,a)=\prod_{u=1}^K P_u(s_u'\mid s_u,a_u)
\]
to hold exactly. The interaction is induced by the Tanner graph itself through shared VNs and cycles.
\end{frame}

\begin{frame}{Step 6: Where the direct restless-bandit mapping fails}
The problem is also not a textbook restless bandit.

\vspace{0.4em}
In a classical restless bandit, each arm has its own active and passive transition kernels:
\[
P_u(s_u'\mid s_u, a_u=1),
\qquad
P_u(s_u'\mid s_u, a_u=0),
\]
independent of the actions taken on other arms.

\vspace{0.4em}
In clustered LDPC scheduling, arm $v$ can be affected by the decision taken on arm $u \neq v$. Therefore the arms are \emph{interacting}. This breaks the exact assumptions behind a direct Whittle-index formulation.
\end{frame}

\begin{frame}{Intermediate conclusion}
For LDPC scheduling with cluster interactions:
\begin{itemize}
  \item it \textbf{does} resemble a weakly coupled MDP in the sense of a shared activation constraint,
  \item it \textbf{does} resemble a restless bandit in the sense of arm-like subsystems that evolve over time,
  \item but it is \textbf{not exactly} either model because the transitions are graph-coupled.
\end{itemize}
The best exact viewpoint is therefore a \textbf{sparsely dynamically coupled factored MDP}.
\end{frame}

\begin{frame}{Step 7: Why RELDEC is still only partially using global information}
RELDEC defines a cluster state from local hard-decision outputs around that cluster and selects actions from local state-dependent action values.

\vspace{0.4em}
This keeps learning feasible, but it does not explicitly encode the current pressure exerted by the rest of the graph on the selected cluster. In other words, the scheduler uses a local state but not the full joint state.

\vspace{0.4em}
That is precisely where an augmented-state formulation can help.
\end{frame}

\begin{frame}{Step 8: State augmentation idea}
For each cluster $u$, define an augmented state
\[
\tilde s_u = (s_u, \phi_u),
\]
where $s_u$ is the local cluster state and $\phi_u$ is a low-dimensional summary of how the rest of the graph currently influences cluster $u$.

\vspace{0.4em}
A generic form is
\[
\phi_u = \sum_{v\neq u} \lambda_{vu} f(s_v),
\]
where $\lambda_{vu}$ measures structural coupling between clusters $v$ and $u$, and $f(s_v)$ extracts a small feature from cluster $v$ such as residual mass, unsatisfied-check fraction, or local error proxy.
\end{frame}

\begin{frame}{Step 9: Why the augmentation is feasible}
The full joint state $s=(s_1,\dots,s_K)$ is usually too large. The augmented state keeps only:
\begin{itemize}
  \item local information $s_u$, and
  \item a compressed coupling statistic $\phi_u$.
\end{itemize}
So the scheduler for cluster $u$ can use a local score such as
\[
g_u(\tilde s_u)=g_u(s_u,\phi_u),
\]
and then activate the top-scoring cluster(s). This preserves the cheap action-selection structure of bandit-style scheduling while injecting global context.

\vspace{0.4em}
Conceptually, this is closer to sparse-interaction factored control, value decomposition, and aggregate-state approximations than to a strict Whittle index.
\end{frame}

\begin{frame}{Step 10: Can we still call it a WC-MDP or a restless bandit?}
\begin{center}
\begin{tabularx}{\textwidth}{>{\raggedright\arraybackslash}p{3.2cm} >{\raggedright\arraybackslash}p{2.6cm} X}
\toprule
Model & Exact fit? & Comment \\
\midrule
Classical WC-MDP & No & Shared budget fits, but transition factorization fails because cluster actions affect neighboring cluster states. \\
Classical restless bandit & No & Arm evolution is not independent under active/passive modes; arms interact through Tanner-graph coupling. \\
Approximate WC-MDP / approximate RMAB & Yes, as approximation & Reasonable when clustering makes cross-cluster interaction weak enough to be treated perturbatively. \\
Dynamically coupled factored MDP & Yes, best exact label & Captures local subsystems with sparse graph-induced transition dependence. \\
\bottomrule
\end{tabularx}
\end{center}
\end{frame}

\begin{frame}{Step 11: Closest existing concepts in the literature}
The augmentation idea is directly related to several existing directions:
\begin{itemize}
  \item \textbf{Factored MDPs}: global state decomposes into interacting local components.
  \item \textbf{Dynamically coupled systems}: local next-states depend on neighboring states/actions.
  \item \textbf{Approximate weak coupling}: use local control plus correction terms when coupling is small.
  \item \textbf{Interacting/networked restless bandits}: bandit-like activation under sparse arm interactions.
  \item \textbf{Aggregate-state / mean-field summaries}: replace the full rest-of-system state by a compressed statistic.
  \item \textbf{Value decomposition / sparse-interaction RL}: local scores with limited interaction terms.
\end{itemize}
\end{frame}

\begin{frame}{Final takeaway}
If we build the argument carefully, the clean academic statement is:

\vspace{0.5em}
\begin{block}{Recommended formulation}
Clustered LDPC scheduling is best modeled as a \textbf{sparsely dynamically coupled factored MDP} with a scheduling-budget constraint. It is \textbf{not exactly} a classical weakly coupled MDP or a classical restless bandit, although both provide useful approximation languages when cross-cluster interactions are weak.
\end{block}

\vspace{0.5em}
The proposed augmentation
\[
\tilde s_u=(s_u,\phi_u)
\]
is then a principled compromise: it preserves local decision structure while importing a compact summary of global coupling.
\end{frame}

\end{document}